Agent-Based Modeling (ABM) was a powerful computational method for simulating the actions and interactions of autonomous agents within a defined environment. Serving as a ``virtual laboratory'' \citep{deSouza2021}, ABM was uniquely suited for modeling complex socio-environmental systems (SES) where macro-level outcomes, such as community compliance, emerged from micro-level behaviors and interactions \citep{Brugiere2022}. This capacity was paramount for analyzing municipal solid waste (MSW), where system-wide compliance results directly from the cumulative decisions made at the household level \citep{Fontaine2024}.

The necessity of ABM in this research was rooted in its ability to model heterogeneity and adaptive behavior. While other methodologies like System Dynamics (SD) were effective for analyzing aggregate stocks and flows of plastic waste \citep{Dhanshyam2021}, and systematic reviews linked macro-population growth to generation rates \citep{Eltanal2025}, these approaches failed to capture individual decision-making. Unlike traditional system dynamics models, ABM captured the diversity of households, representing them as agents with distinct socio-demographic factors and psychological profiles, as defined by the Theory of Planned Behavior (TPB) \citep{Fontaine2024}. This allowed the simulation of non-linear and adaptable responses to policy changes---a critical feature since household behaviors were not static. Furthermore, ABM's capacity to integrate advanced computational techniques, such as incorporating machine learning (ML) classifiers into agent decision logic, enhanced behavioral realism beyond static heuristics and improved the accuracy of predicted policy outcomes \citep{Jimenez2025}.
